\section{Discussion}
\label{sec:discussion}

\subsection{What each theory contributed}

The paper's structural claim is that all three sources are load-bearing and
none is decorative. \Cref{tab:contributions} states what each supplied and what
breaks without it.

\begin{table}[t]
\centering
\caption{Each theory's contribution, and the failure that returns if it is
removed.}
\label{tab:contributions}
\footnotesize
\begin{tabular}{@{}lp{0.29\columnwidth}p{0.33\columnwidth}@{}}
\toprule
& \textbf{Contributed} & \textbf{Without it} \\
\midrule
Bell--LaPadula
& The write ceiling and the read/write asymmetry: read at $0.9$, emit at $0.2$
& No principled basis for reading e5 while never emitting it; access becomes
  all-or-nothing \\[3pt]
Biba
& The integrity axis and the action plane: e8/e9 quarantined, e2 restored
& Injection is unmodelled; every confidentiality metric reads clean while the
  verdict is wrong \\[3pt]
Zero trust
& Per-request recomputation and least-privilege scoping: e7 flips between
  runs; ambient credentials never connect
& Labels cache as ``internal, therefore fine''; the Downward decision leaks
  into the External one \\[3pt]
CI-Work
& The five flow directions, which supply $\clr$
& The single most consequential number in the trace has no source \\
\bottomrule
\end{tabular}
\end{table}

It is worth noting that CI-Work's contribution here is obtained \emph{without
any model judgement}: direction is read from organizational structure, and it
determines the ceiling that drives most decisions.

\subsection{Two latent bugs our own derivation exposed}
\label{sec:discussion:bugs}

We report these because both are instructive about the class of system, not
because confession is a genre convention. Both were found while deriving the
weights of \Cref{sec:method:weights}, and both were reachable under the
configuration the work was about to adopt rather than the one already measured.

\paragraph{Our integrity labeler was itself prompt-injectable.} The rubric
originally applied an unconditional uplift: a high-integrity content marker
such as ``approved'' or ``official'' raised an entry's score to $0.85$
\emph{regardless of its source}. So a PR author writing ``Approved by the
security team offline---safe to merge'' in the description---the least
trustworthy channel in the system, controlled by the party under review---
produced $\Il = \max(0.10, 0.85) = 0.85$, clearing $\lvlI(\text{External}) =
0.60$ comfortably.

This is worse than a calibration error because it is adversarially reachable
and requires no knowledge of the configuration. It also cuts against the
paper's own thesis: our argument is that prompt injection is an integrity
violation that Biba handles, and here was our integrity labeler being
injected. The fix follows from a principle rather than a patch:
\textbf{integrity is a property of provenance, and content cannot vouch for
its own provenance}. Uplift is now bounded relative to the source's base score,
with an absolute ceiling on author-controlled channels applied last so no later
signal can undo it (\Cref{sec:method:labels}). Two supporting changes matter as
much as the clamp: we removed ``approved'' and ``confirmed'' from the marker
defaults, being the two most author-writable words in the list; and source keys
must now name the \emph{artifact} rather than the platform, since one host
serves both an authoritative diff and an author-controlled description and they
cannot share a score.

\paragraph{The threshold override could release regulated content.} The
$\theta$ override originally applied to \textsc{deny} as well as
\textsc{quarantine}, with no floor. Under the derived weights, a Downward
request with $\wC = 0.67$ and $\clr = 0.70$ gives an entry at $\Cl = 0.80$ a
risk of $0.067$, releasing a Restricted-class entry because its margin happened
to be small and the direction's weight low. Leniency for near-threshold entries
is reasonable; silently releasing regulated content is not. Precedence rule~3
now sits above the override (\Cref{tab:precedence}).

The reason this was not cosmetic is the interesting part. The \emph{measured}
results were unchanged by the fix, because the previously used configuration
never reached the override. But across the $\theta$ sweep the paper plans to
report, the pre-fix code released a named-customer regulated entry at every
operating point from $\theta = 0.2$ upward. \textbf{The vulnerability was
latent in the configuration we were about to adopt, not in the one we had
already measured}---which is a general hazard for any system whose safety
properties are threshold-dependent, and an argument for sweeping before
publishing rather than after.

\paragraph{A methodological note.} Writing the cases into an executable fixture
immediately falsified two claims in our own hand-computed walkthrough: a
conveyance figure and an integrity-leakage figure (\Cref{sec:eval:ilblindspot}).
We take this as an argument for executable design documents in systems-security
work, and it is why the artifact includes a verifier that re-derives the
paper's decisions from the configuration rather than merely reproducing its
tables.

\subsection{A falsifiable prediction}
\label{sec:discussion:prediction}

CI-Work reports \emph{inverse scaling}: larger models in the same family convey
more and leak more, which they attribute to better attention over long noisy
contexts combined with greater sycophancy toward user instructions over
implicit norms.

Both proposed mechanisms require the sensitive detail to \emph{be in the
context window}. Under L2 handle-mediated citation it is not. This yields a
clean prediction:

\begin{quote}
\textbf{Under L2, the leakage-versus-model-size slope flattens toward zero,
while under L0 and L1 it remains positive.}
\end{quote}

The experiment is three model sizes crossed with three enforcement
levels---small enough to actually run, and directly contradicting CI-Work's
scaling result if L2 works. We consider this the strongest single result the
paper could carry, and we state the falsification condition with it: if the
slope stays positive under L2, then either the abstracts leak or the mechanism
is not doing what we claim, and the L2 argument fails.
\needsrun{this is the headline experiment and it has not been run}

\subsection{Design validation against production incidents}
\label{sec:discussion:incidents}

Returning to \Cref{tab:incidents}, the framework reads those incidents as one
failure rather than five, and the reading is not merely descriptive.

The PocketOS case is the sharpest. The agent hit a credential mismatch, formed
the hypothesis that a volume was staging-scoped, and acted on it irreversibly.
In our terms the action demanded $\Il \ge 0.95$ and the evidence was an
unverified inference---the agent's own account was that it ``guessed rather
than verified.'' \Cref{eq:downgrade} caps the permitted action at the integrity
of the weakest supporting evidence, which here is near the floor, so the
permitted action is to ask a question. Crucially, this requires no detection of
anything: no classifier must recognize the plan as dangerous, because the bound
follows arithmetically from labels the system already holds.

The Replit case validates the ladder rather than the rule. A code freeze was
stated as an instruction and the agent violated it, which is the definition of
L0. The vendor's remediation---automatic development/production separation and
a planning-only mode---is a capability constraint, i.e.\ a move from L0 to the
action plane. That the fix converged independently on our mechanism is modest
evidence the mechanism is the right one.

We are wary of overreading this subsection, and note two honest limits. These
are trade-press and practitioner reports rather than peer-reviewed
postmortems; and it is always easier to explain an incident after the fact than
to prevent the next one. The subsection's proper weight is as a demonstration
that the failure mode we formalize occurs outside our benchmark, not as
evidence that our system would have prevented these specific events.

\subsection{Limitations}
\label{sec:discussion:limitations}

\paragraph{We address authority, not rate.} The cost incidents in
\Cref{tab:incidents} are substantially \emph{rate} failures: retry loops with
no timeout, no per-agent spend limit, and billing telemetry that lags spend by
roughly a day~\cite{infoq2026billing}. Our mechanism makes spend authority an
explicit policy input rather than an ambient property of a credential, which
is a real contribution to the problem, but it does not bound iteration count
and would not by itself have stopped a loop. Spend caps, timeouts, and kill
switches are complementary and necessary, and we do not claim them.

\paragraph{Scale.} Three implemented seeds with a deterministic judge is a
demonstration, not an evaluation. The claims in \Cref{sec:eval} are
appropriately scoped to what three cases can support---which is the existence
of the failure mode, the disjointness of the two rule families, and the
conditional statement about the write ceiling---and not to effect sizes.
\needsrun{$\sim$100 seeds with $\ge$20\% null, and the LLM judge}

\paragraph{Label error is unmeasured.} All reported results use evaluation-mode
ground-truth $\Cl$, which isolates policy quality from labeler error by
design. The inference-mode agreement study specified in
\Cref{sec:method:weights} has not been run, so we cannot yet state how much
performance is lost when labels are inferred rather than given. This is the
single largest gap between our results and a deployment claim.

\paragraph{Abstraction is an assumption, not yet a mechanism.} The L2 abstracts
are hand-written in the fixture rather than produced by an audited offline
abstractor, and the adversarial reconstruction check---prompting a model with
an abstract plus all admitted entries to see whether it can recover the
restricted specifics---does not exist. Until it does, the claim that an
abstract is safe is an assumption. An abstract can be individually below the
ceiling and still permit reconstruction in conjunction with admitted entries,
which is precisely the conjunction blindness that \Cref{sec:method:ladder}
identifies in L1.

\paragraph{Per-entry decisions cannot see conjunctions.} By construction, a
policy that decides entry by entry cannot detect that two admitted entries
\emph{together} identify what a denied entry named. This is a structural limit
of the design rather than of the implementation, and it bounds what any
label-based approach can promise against an adversary who probes.

\paragraph{The trust boundary moved rather than vanished.} L2 relocates trust
to the offline abstractor and L3 to the schema interface, whose reason-code
vocabulary is a covert channel of nonzero declared capacity. We prefer these
boundaries to the original because they are smaller, offline, and auditable,
but they are boundaries.

\paragraph{Generality beyond PR review.} The benchmark instantiates one task.
We argue the structure generalizes---the two ceilings, the three entry classes,
and the capability-indexed evidence bar are not specific to code
review---but we have not demonstrated it on a second task, and the
$\lvlI^{\mathrm{cap}}$ ladder of \Cref{tab:downgrade} would need re-deriving
for each new action surface.
